% --- Archon physics formalization source begin ---
% archon:physics
% archon:covers PhyXMiniProblems/problem_phyx_mini_0067.lean
% archon:source-report reports/phyx_mini/problem_phyx_mini_0067.source.json

\chapter{Physics problem phyx\_mini\_0067}
\label{ch:PhyXMiniProblems_problem_phyx_mini_0067}

\paragraph{Problem source.}
\#\# Question

Find the focal length of the meniscus polystyrene plastic lens

\#\# Answer choices

- A: A: 24cm

- B: B: 56cm

- C: C: 20cm

- D: D: 28cm

\#\# Figure caption (auxiliary; use the image as primary evidence)

The image depicts a meniscus lens, which is a type of lens with concave and convex surfaces. The lens is positioned vertically, with its curved surfaces illustrated in blue. Two lines with arrows indicate distances measured from a straight horizontal line passing through the center of the lens. One line is labeled "30 cm," representing a distance from the lens to a point on the right. The other line is labeled "40 cm," indicating a further distance to another point on the same side. The text "Meniscus lens" is written below the image. The points seem to correspond to focal points or object/image positions related to the lens.

\#\# Dataset metadata

Geometrical Optics; Physical Model Grounding Reasoning, Multi-Formula Reasoning

\paragraph{Recorded answer/context.}
C: C: 20cm

\paragraph{Figure/image path.}
/root/proposal\_for\_physic/science-mango/phyx\_mini\_run/phyx\_data/test\_image/67.png

\paragraph{Formalization target.}
create a compiling Lean file with sorry bodies at `PhyXMiniProblems/problem\_phyx\_mini\_0067.lean`. The Lean declarations must preserve the physical quantities, dimensions or dimensional roles, figure labels, governing-law hypotheses, and final relation expressed by this problem.
Use Mathlib/Physlib names found through LeanExplore where available. If a physics API is missing, introduce faithful local abstractions rather than scalar placeholder aliases.

\begin{theorem}[Physics formalization target]
\label{thm:physics:phyx_mini_0067:target}
\lean{PhyXMiniProblems.ProblemPhyXMini0067.problem_phyx_mini_0067}
\uses{def:physics:phyx-mini-0067:phyxminiproblems-problemphyxmini0067-lengthreadout, def:physics:phyx-mini-0067:phyxminiproblems-problemphyxmini0067-meniscuslenssetup, def:physics:phyx-mini-0067:phyxminiproblems-problemphyxmini0067-matchesprimaryfigure, def:physics:phyx-mini-0067:phyxminiproblems-problemphyxmini0067-matchespolystyrenethinlensmodel, def:physics:phyx-mini-0067:phyxminiproblems-problemphyxmini0067-hasphysicalconfiguration, def:physics:phyx-mini-0067:phyxminiproblems-problemphyxmini0067-satisfiesthinlensmakerequation, def:physics:phyx-mini-0067:phyxminiproblems-problemphyxmini0067-answerchoice, def:physics:phyx-mini-0067:phyxminiproblems-problemphyxmini0067-matchesdisplayedfocallength, def:physics:phyx-mini-0067:phyxminiproblems-problemphyxmini0067-recordedanswerclaim}
The assigned autoformalize agent should translate this physics problem into Lean declarations in the covered file, with theorem and lemma proof bodies written as `by sorry`.
\end{theorem}
\begin{proof}
This is an autoformalization task, not a proof task. Produce statements that can later be proved by the physics prover without weakening the physical model.
\end{proof}

% --- Archon declaration topology begin ---
\section{Declaration topology}
\begin{definition}[Optical Length]
  \label{def:physics:phyx-mini-0067:phyxminiproblems-problemphyxmini0067-opticallength}
  \lean{PhyXMiniProblems.ProblemPhyXMini0067.OpticalLength}
  A signed physical optical length, independent of a choice of units
\end{definition}
\begin{proof}
  This is the declared model definition of Optical Length.
\end{proof}
\begin{definition}[length Readout]
  \label{def:physics:phyx-mini-0067:phyxminiproblems-problemphyxmini0067-lengthreadout}
  \lean{PhyXMiniProblems.ProblemPhyXMini0067.lengthReadout}
  \uses{def:physics:phyx-mini-0067:phyxminiproblems-problemphyxmini0067-opticallength}
  The signed scalar readout of an optical length in the selected unit
\end{definition}
\begin{proof}
  This is the declared model definition of length Readout.
\end{proof}
\begin{definition}[Optical Medium]
  \label{def:physics:phyx-mini-0067:phyxminiproblems-problemphyxmini0067-opticalmedium}
  \lean{PhyXMiniProblems.ProblemPhyXMini0067.OpticalMedium}
  The homogeneous optical media relevant to the lensmaker model
\end{definition}
\begin{proof}
  This is the declared model definition of Optical Medium.
\end{proof}
\begin{definition}[Lens Profile]
  \label{def:physics:phyx-mini-0067:phyxminiproblems-problemphyxmini0067-lensprofile}
  \lean{PhyXMiniProblems.ProblemPhyXMini0067.LensProfile}
  The cross-sectional lens profile explicitly named in the problem
\end{definition}
\begin{proof}
  This is the declared model definition of Lens Profile.
\end{proof}
\begin{definition}[Optical Approximation]
  \label{def:physics:phyx-mini-0067:phyxminiproblems-problemphyxmini0067-opticalapproximation}
  \lean{PhyXMiniProblems.ProblemPhyXMini0067.OpticalApproximation}
  The optical approximation used by the governing lensmaker equation
\end{definition}
\begin{proof}
  This is the declared model definition of Optical Approximation.
\end{proof}
\begin{definition}[Lens Surface]
  \label{def:physics:phyx-mini-0067:phyxminiproblems-problemphyxmini0067-lenssurface}
  \lean{PhyXMiniProblems.ProblemPhyXMini0067.LensSurface}
  The two refracting surfaces in left-to-right propagation order
\end{definition}
\begin{proof}
  This is the declared model definition of Lens Surface.
\end{proof}
\begin{definition}[Axis Point]
  \label{def:physics:phyx-mini-0067:phyxminiproblems-problemphyxmini0067-axispoint}
  \lean{PhyXMiniProblems.ProblemPhyXMini0067.AxisPoint}
  Axial points determined by the two spherical-surface constructions
\end{definition}
\begin{proof}
  This is the declared model definition of Axis Point.
\end{proof}
\begin{definition}[vertex]
  \label{def:physics:phyx-mini-0067:phyxminiproblems-problemphyxmini0067-lenssurface-vertex}
  \lean{PhyXMiniProblems.ProblemPhyXMini0067.LensSurface.vertex}
  \uses{def:physics:phyx-mini-0067:phyxminiproblems-problemphyxmini0067-lenssurface, def:physics:phyx-mini-0067:phyxminiproblems-problemphyxmini0067-axispoint}
  The axial vertex belonging to a refracting surface
\end{definition}
\begin{proof}
  This is the declared model definition of vertex.
\end{proof}
\begin{definition}[center Of Curvature]
  \label{def:physics:phyx-mini-0067:phyxminiproblems-problemphyxmini0067-lenssurface-centerofcurvature}
  \lean{PhyXMiniProblems.ProblemPhyXMini0067.LensSurface.centerOfCurvature}
  \uses{def:physics:phyx-mini-0067:phyxminiproblems-problemphyxmini0067-lenssurface, def:physics:phyx-mini-0067:phyxminiproblems-problemphyxmini0067-axispoint}
  The axial center of curvature belonging to a refracting surface
\end{definition}
\begin{proof}
  This is the declared model definition of center Of Curvature.
\end{proof}
\begin{definition}[Meniscus Lens Setup]
  \label{def:physics:phyx-mini-0067:phyxminiproblems-problemphyxmini0067-meniscuslenssetup}
  \lean{PhyXMiniProblems.ProblemPhyXMini0067.MeniscusLensSetup}
  \uses{def:physics:phyx-mini-0067:phyxminiproblems-problemphyxmini0067-opticallength, def:physics:phyx-mini-0067:phyxminiproblems-problemphyxmini0067-opticalmedium, def:physics:phyx-mini-0067:phyxminiproblems-problemphyxmini0067-lensprofile, def:physics:phyx-mini-0067:phyxminiproblems-problemphyxmini0067-opticalapproximation, def:physics:phyx-mini-0067:phyxminiproblems-problemphyxmini0067-axispoint}
  Physical quantities attached to the polystyrene meniscus-lens diagram. axialPosition is a signed coordinate on the displayed optical axis; focalLength is the signed effective focal length of the complete lens
\end{definition}
\begin{proof}
  This is the declared model definition of Meniscus Lens Setup.
\end{proof}
\begin{definition}[axial Separation Readout]
  \label{def:physics:phyx-mini-0067:phyxminiproblems-problemphyxmini0067-axialseparationreadout}
  \lean{PhyXMiniProblems.ProblemPhyXMini0067.axialSeparationReadout}
  \uses{def:physics:phyx-mini-0067:phyxminiproblems-problemphyxmini0067-lengthreadout, def:physics:phyx-mini-0067:phyxminiproblems-problemphyxmini0067-axispoint, def:physics:phyx-mini-0067:phyxminiproblems-problemphyxmini0067-meniscuslenssetup}
  Signed axial separation from start to finish
\end{definition}
\begin{proof}
  This is the declared model definition of axial Separation Readout.
\end{proof}
\begin{definition}[signed Surface Radius Readout]
  \label{def:physics:phyx-mini-0067:phyxminiproblems-problemphyxmini0067-signedsurfaceradiusreadout}
  \lean{PhyXMiniProblems.ProblemPhyXMini0067.signedSurfaceRadiusReadout}
  \uses{def:physics:phyx-mini-0067:phyxminiproblems-problemphyxmini0067-lenssurface, def:physics:phyx-mini-0067:phyxminiproblems-problemphyxmini0067-meniscuslenssetup, def:physics:phyx-mini-0067:phyxminiproblems-problemphyxmini0067-axialseparationreadout}
  The signed spherical radius of a surface. It is positive when its center of curvature is to the right of its vertex on the oriented optical axis
\end{definition}
\begin{proof}
  This is the declared model definition of signed Surface Radius Readout.
\end{proof}
\begin{definition}[Matches Primary Figure]
  \label{def:physics:phyx-mini-0067:phyxminiproblems-problemphyxmini0067-matchesprimaryfigure}
  \lean{PhyXMiniProblems.ProblemPhyXMini0067.MatchesPrimaryFigure}
  \uses{def:physics:phyx-mini-0067:phyxminiproblems-problemphyxmini0067-lengthreadout, def:physics:phyx-mini-0067:phyxminiproblems-problemphyxmini0067-meniscuslenssetup, def:physics:phyx-mini-0067:phyxminiproblems-problemphyxmini0067-signedsurfaceradiusreadout}
  Data read directly from the primary figure: a meniscus profile, radius labels 30 cm and 40 cm, and the depicted left-to-right ordering of the vertices and centers. No focal-length value occurs in this predicate
\end{definition}
\begin{proof}
  This is the declared model definition of Matches Primary Figure.
\end{proof}
\begin{definition}[Matches Polystyrene Thin Lens Model]
  \label{def:physics:phyx-mini-0067:phyxminiproblems-problemphyxmini0067-matchespolystyrenethinlensmodel}
  \lean{PhyXMiniProblems.ProblemPhyXMini0067.MatchesPolystyreneThinLensModel}
  \uses{def:physics:phyx-mini-0067:phyxminiproblems-problemphyxmini0067-meniscuslenssetup}
  Material-table and modeling data implicit in “polystyrene plastic lens”: the lens is in air, the idealized refractive indices are 1.6 and 1, and the thin paraxial approximation is selected. These are input data rather than a statement of the requested focal length
\end{definition}
\begin{proof}
  This is the declared model definition of Matches Polystyrene Thin Lens Model.
\end{proof}
\begin{definition}[Has Physical Configuration]
  \label{def:physics:phyx-mini-0067:phyxminiproblems-problemphyxmini0067-hasphysicalconfiguration}
  \lean{PhyXMiniProblems.ProblemPhyXMini0067.HasPhysicalConfiguration}
  \uses{def:physics:phyx-mini-0067:phyxminiproblems-problemphyxmini0067-lengthreadout, def:physics:phyx-mini-0067:phyxminiproblems-problemphyxmini0067-meniscuslenssetup, def:physics:phyx-mini-0067:phyxminiproblems-problemphyxmini0067-signedsurfaceradiusreadout}
  Positivity and nondegeneracy conditions for the depicted optical branch
\end{definition}
\begin{proof}
  This is the declared model definition of Has Physical Configuration.
\end{proof}
\begin{definition}[Satisfies Thin Lensmaker Equation]
  \label{def:physics:phyx-mini-0067:phyxminiproblems-problemphyxmini0067-satisfiesthinlensmakerequation}
  \lean{PhyXMiniProblems.ProblemPhyXMini0067.SatisfiesThinLensmakerEquation}
  \uses{def:physics:phyx-mini-0067:phyxminiproblems-problemphyxmini0067-lengthreadout, def:physics:phyx-mini-0067:phyxminiproblems-problemphyxmini0067-meniscuslenssetup, def:physics:phyx-mini-0067:phyxminiproblems-problemphyxmini0067-signedsurfaceradiusreadout}
  The thin-lens lensmaker equation in a common surrounding medium, 1 / f = (n\_lens / n\_medium - 1) (1 / R\_front - 1 / R\_back). It is stated for every length unit, so both sides transform as inverse length. This governing law contains no numerical focal-length answer
\end{definition}
\begin{proof}
  This is the declared model definition of Satisfies Thin Lensmaker Equation.
\end{proof}
\begin{definition}[Answer Choice]
  \label{def:physics:phyx-mini-0067:phyxminiproblems-problemphyxmini0067-answerchoice}
  \lean{PhyXMiniProblems.ProblemPhyXMini0067.AnswerChoice}
  Labels of the four answer choices printed with the problem
\end{definition}
\begin{proof}
  This is the declared model definition of Answer Choice.
\end{proof}
\begin{definition}[displayed Focal Length In Centimeters]
  \label{def:physics:phyx-mini-0067:phyxminiproblems-problemphyxmini0067-displayedfocallengthincentimeters}
  \lean{PhyXMiniProblems.ProblemPhyXMini0067.displayedFocalLengthInCentimeters}
  \uses{def:physics:phyx-mini-0067:phyxminiproblems-problemphyxmini0067-answerchoice}
  The whole-centimeter focal length printed beside each answer choice
\end{definition}
\begin{proof}
  This is the declared model definition of displayed Focal Length In Centimeters.
\end{proof}
\begin{definition}[recorded Answer Choice]
  \label{def:physics:phyx-mini-0067:phyxminiproblems-problemphyxmini0067-recordedanswerchoice}
  \lean{PhyXMiniProblems.ProblemPhyXMini0067.recordedAnswerChoice}
  \uses{def:physics:phyx-mini-0067:phyxminiproblems-problemphyxmini0067-answerchoice}
  The answer label recorded by the dataset
\end{definition}
\begin{proof}
  This is the declared model definition of recorded Answer Choice.
\end{proof}
\begin{definition}[Matches Displayed Focal Length]
  \label{def:physics:phyx-mini-0067:phyxminiproblems-problemphyxmini0067-matchesdisplayedfocallength}
  \lean{PhyXMiniProblems.ProblemPhyXMini0067.MatchesDisplayedFocalLength}
  \uses{def:physics:phyx-mini-0067:phyxminiproblems-problemphyxmini0067-lengthreadout, def:physics:phyx-mini-0067:phyxminiproblems-problemphyxmini0067-meniscuslenssetup, def:physics:phyx-mini-0067:phyxminiproblems-problemphyxmini0067-answerchoice, def:physics:phyx-mini-0067:phyxminiproblems-problemphyxmini0067-displayedfocallengthincentimeters}
  Exact agreement between the modeled focal length and a displayed choice
\end{definition}
\begin{proof}
  This is the declared model definition of Matches Displayed Focal Length.
\end{proof}
\begin{definition}[Recorded Answer Claim]
  \label{def:physics:phyx-mini-0067:phyxminiproblems-problemphyxmini0067-recordedanswerclaim}
  \lean{PhyXMiniProblems.ProblemPhyXMini0067.RecordedAnswerClaim}
  \uses{def:physics:phyx-mini-0067:phyxminiproblems-problemphyxmini0067-meniscuslenssetup, def:physics:phyx-mini-0067:phyxminiproblems-problemphyxmini0067-recordedanswerchoice, def:physics:phyx-mini-0067:phyxminiproblems-problemphyxmini0067-matchesdisplayedfocallength}
  The source's recorded claim that choice C, 20 cm, is the answer
\end{definition}
\begin{proof}
  This is the declared model definition of Recorded Answer Claim.
\end{proof}
\begin{lemma}[focal Length In Centimeters eq two Hundred]
  \label{lem:physics:phyx-mini-0067:phyxminiproblems-problemphyxmini0067-focallengthincentimeters-eq-twohundred}
  \lean{PhyXMiniProblems.ProblemPhyXMini0067.focalLengthInCentimeters_eq_twoHundred}
  \uses{def:physics:phyx-mini-0067:phyxminiproblems-problemphyxmini0067-lengthreadout, def:physics:phyx-mini-0067:phyxminiproblems-problemphyxmini0067-meniscuslenssetup, def:physics:phyx-mini-0067:phyxminiproblems-problemphyxmini0067-matchesprimaryfigure, def:physics:phyx-mini-0067:phyxminiproblems-problemphyxmini0067-matchespolystyrenethinlensmodel, def:physics:phyx-mini-0067:phyxminiproblems-problemphyxmini0067-hasphysicalconfiguration, def:physics:phyx-mini-0067:phyxminiproblems-problemphyxmini0067-satisfiesthinlensmakerequation}
  With the image's 30 cm and 40 cm signed radii and the idealized polystyrene index 1.6, the thin-lens lensmaker equation gives f = 200 cm
\end{lemma}
\begin{proof}
  The conclusion follows from the governing relations and figure readouts named in the statement of focal Length In Centimeters eq two Hundred.
\end{proof}
% --- Archon declaration topology end ---
% --- Archon physics formalization source end ---
